\begin{definition}[Ambrosio--Kirchheim mass-supremum formula, hard direction, finite form: the
  statement, taken as an explicit hypothesis]
  Let $E$ be a complete, separable metric space equipped with a compatible Borel $\sigma$-algebra.
  We write $\mathrm{MassSupFormulaStatement}(E)$ for the following proposition about $E$:
  \begin{quote}
    for every $T \in \mathcal{M}_1(E)$ (\noderef{MetricCurrent1}) and every $c \in \mathbb{R}$
    with $c \ge 0$: if for every $n \in \mathbb{N}$ and every finite family $(f_p,\pi_p)_{p<n}
    \in \mathcal{D}^1(E)$ (\noderef{Form1}) with $\mathrm{Lip}(\pi_p) \le 1$ for all $p<n$ and
    $\sum_{p<n} |f_p(x)| \le 1$ for every $x \in E$ one has
    \[
      \sum_{p<n} |T(f_p,\pi_p)| \le c
      ,
    \]
    then $\mathbb{M}(T) \le c$, i.e.\ $\mathrm{massOfFunctional}(T) \le
    \mathrm{ENNReal.ofReal}(c)$ (\noderef{massOfFunctional}).
  \end{quote}

  \textbf{Why this node is a statement rather than a theorem.} This node is a verbatim external
  literature citation, \cite{AmbrosioKirchheim}*{Proposition 2.7}, which the tablet deliberately
  does not reprove. Rather than record it as a theorem with an unproved body -- which would make
  every downstream result depend on \texttt{sorryAx} -- it is recorded here as a named proposition
  about $E$, and the nodes that consume it (\noderef{BBAssertion}, and through it
  \noderef{ApproximationMetric}) carry $\mathrm{MassSupFormulaStatement}(E)$ as an explicit
  hypothesis, discharged by whoever invokes them. The remainder of this node justifies that the
  proposition displayed above is exactly the cited result, and records the reductions used to put
  it in this form.

  \textbf{Hypotheses supplied.} $E$ is complete and separable by the ambient
  \texttt{[CompleteSpace E] [TopologicalSpace.SeparableSpace E]} instances, matching the paper's
  standing assumption at paper.tex line 1096 (``Throughout this paper $E$ \dots\ denotes a
  complete, separable, geodesic metric space'') under which \cite{AmbrosioKirchheim}*{Proposition
  2.7} is invoked at paper.tex 1497--1501. We do \emph{not} add a geodesic hypothesis (e.g.\ a
  geodesic-pairing structure on $E$) here: Proposition 2.7 is a statement purely about the mass
  $\mathbb{M}(T)$ of a $1$-current $T \in \mathcal{M}_1(E)$ and its representation as a supremum
  over $1$-forms, with no reference to geodesics anywhere in its statement or in
  \cite{AmbrosioKirchheim}'s proof of it; the paper's standing assumption bundles geodesy into its
  blanket hypothesis on $E$ only because $E$'s later role as the ambient space for curve
  constructions (Section~2 onward) needs it, not because Proposition 2.7 does. (The tablet's other
  citation-only external-input node for this same standing hypothesis, ``PaoliniStepanovExists'',
  takes the same line: it carries completeness and separability but not geodesy.)

  We also do not add continuity of $T$ in the $f$-slot as a hypothesis, even though
  \cite{AmbrosioKirchheim}*{Definition 3.1} of a metric $1$-current includes it: as
  \noderef{MetricCurrent1}'s own definition records, this continuity is not axiomatized there
  because paper.tex line 1393 derives it from \texttt{finiteMass} together with Lebesgue's
  dominated convergence theorem (a derivation the tablet supplies elsewhere, in a node named
  ``CurrentContinuousF''). Since $T : \mathrm{MetricCurrent1}\ E$ already carries
  \texttt{finiteMass}, this continuity is available whenever it is needed and need not be repeated
  as a hypothesis of this node.

  \textbf{The external theorem, verbatim.} \cite{AmbrosioKirchheim}*{Proposition 2.7} (paper.tex
  1497--1501) asserts that for a $1$-current $T$,
  \[
    \mathbb{M}(T) = \sup \sum_{p=1}^\infty |T(f_p,\pi_p)|
    ,
  \]
  where the supremum is over sequences $f_p \in \mathrm{Lip}_b(E)$, $\pi_p \in \mathrm{Lip}(E)$
  with $\mathrm{Lip}(\pi_p) \le 1$ for all $p$ and $\sum_p |f_p| \le 1$ (i.e.\ $\sum_{p=1}^\infty
  |f_p(x)| \le 1$ for every $x \in E$).

  \textbf{Why only the hard direction is stated, and why this is a citation and not an omission.}
  Writing this equality as two inequalities, $\sup \le \mathbb{M}(T)$ and $\mathbb{M}(T) \le \sup$:
  the first is \emph{free} from the definition of $\mathbb{M}(T) = \mathrm{massOfFunctional}(T)$
  as an infimum over admissible measures (\noderef{massOfFunctional}), and must be (and is)
  re-derived at each use site rather than cited here: for admissible $\mu$ and any such sequence,
  \[
    \sum_{p=1}^\infty |T(f_p,\pi_p)|
    \le \sum_{p=1}^\infty \mathrm{Lip}(\pi_p) \int_E |f_p| \, d\mu
    \le \sum_{p=1}^\infty \int_E |f_p| \, d\mu
    = \int_E \sum_{p=1}^\infty |f_p| \, d\mu
    \le \mu(E)
    ,
  \]
  using $\mathrm{Lip}(\pi_p) \le 1$, monotone convergence to interchange sum and integral, and
  $\sum_p |f_p| \le 1$; taking the infimum over admissible $\mu$ gives $\sup \le \mathbb{M}(T)$
  with no external input. The second direction, $\mathbb{M}(T) \le \sup$, is the content of this
  node. With $\mathbb{M}(T)$ defined as an \emph{infimum} over admissible measures rather than as
  the total mass of a known minimizer, this direction requires \emph{constructing} an admissible
  measure of total mass $\le c$ out of a bound $c$ on the sup -- precisely the mass-measure
  construction that is Ambrosio--Kirchheim's own achievement in proving Proposition 2.7, which this
  tablet deliberately does not reprove (\noderef{massOfFunctional}'s modeling note). It is also not
  recoverable by a soft argument from what the tablet already has: one might hope to obtain an
  admissible measure of mass $\le c$ as a weak-* limit of a sequence of admissible measures with
  masses decreasing to $\inf$, but weak-* compactness of a bounded set of finite Borel measures on
  a general complete separable metric space requires \emph{tightness}, which is not available
  without further hypotheses on $T$ or $E$; so lower semicontinuity of $\mu \mapsto \mu(E)$ under
  such limits cannot be invoked either. Hence $\mathbb{M}(T) \le \sup$ is recorded here as the
  genuine external input, matching the paper's own practice of citing Proposition 2.7 without
  proof.

  \textbf{Why the finite (\texttt{Fin n}-indexed) form is equivalent to, and sufficient in place
  of, the countable form, and why a \texttt{tsum} formulation is avoided.} Since every term
  $|T(f_p,\pi_p)|$ is nonnegative, the infinite series $\sum_{p=1}^\infty |T(f_p,\pi_p)|$ is by
  definition the supremum over $n$ of its partial sums $\sum_{p<n} |T(f_p,\pi_p)|$, so
  \[
    \sup_{(f_p,\pi_p)_{p\ge 0}} \sum_{p=1}^\infty |T(f_p,\pi_p)|
    = \sup_{(f_p,\pi_p)_{p \ge 0}} \ \sup_{n} \sum_{p<n} |T(f_p,\pi_p)|
    = \sup_{n} \ \sup_{(f_p,\pi_p)_{p<n}} \sum_{p<n} |T(f_p,\pi_p)|
    ,
  \]
  where the last equality exchanges the two suprema (valid since both are suprema of the same
  jointly-indexed nonnegative quantity) and drops the dependence of the inner supremum on any
  $\pi_p,f_p$ with $p \ge n$. Concretely: given a sequence $(f_p,\pi_p)_{p\ge 0}$ satisfying the
  two side conditions, its length-$n$ truncation $(f_p,\pi_p)_{p<n}$ satisfies the same two
  conditions restricted to $p<n$ (a sub-sum of a pointwise bound $\le 1$ is itself $\le 1$) and
  realizes the same partial sum. Conversely, given a finite family $(f_p,\pi_p)_{p<n}$ satisfying
  the two side conditions, extend it to a sequence by setting $f_p := 0$, $\pi_p := 0$ for $p \ge
  n$ (valid witnesses: bound $0$, Lipschitz constant $0$); this preserves both side conditions
  ($\mathrm{Lip}(0) = 0 \le 1$; the pointwise sum $\sum_p |f_p(x)|$ is unchanged since the added
  terms are $0$) and contributes $0$ to every partial sum beyond $n$, since $T(0,0) = 0$ by
  \noderef{MetricCurrent1}'s homogeneity in $f$ at scalar $c = 0$ (taking $\omega_c = (0,0)$ and
  $\omega = \omega_c$: the hypotheses of that field hold with $\omega_c.f \equiv 0 = 0 \cdot
  \omega_c.f$, giving $T(0,0) = 0 \cdot T(0,0) = 0$). Hence the two suprema -- over sequences of
  the full series, and over $n$ and $\mathtt{Fin}\ n$-indexed families of the $n$-term partial sum
  -- coincide, and the hypothesis $\mathtt{hsup}$ (a bound $c$ on every such finite partial sum,
  for every $n$) is exactly the statement $\sup \le c$ in Proposition 2.7's formula. We index the
  hypothesis over $\mathtt{Fin}\ n \to \mathrm{Form1}\ E$ rather than as a $\mathtt{tsum}$ over
  $\mathbb{N}$: Lean's $\mathtt{tsum}$ of a non-summable family is $0$ by convention, so a
  $\mathtt{tsum}$-indexed hypothesis $\sum'_p |f_p(x)| \le 1$ would be satisfied vacuously by
  families that do not actually have $\sum_p |f_p(x)| \le 1$ as a genuine bound on partial sums,
  silently weakening the hypothesis; the finite form has no such hazard, and is exactly what a
  countable-indexed bound on partial sums (as produced, for instance, by Fatou's lemma for series
  at the use site of this node) supplies directly.

  \textbf{Conclusion.} By the equivalence above, $\mathtt{hsup}$ gives
  $\sup \sum_{p=1}^\infty |T(f_p,\pi_p)| \le c$ over all admissible sequences $(f_p,\pi_p)$. By
  Proposition 2.7's equality $\mathbb{M}_{\mathrm{AK}}(T) = \sup \sum_{p=1}^\infty
  |T(f_p,\pi_p)|$ (writing $\mathbb{M}_{\mathrm{AK}}(T)$ for the paper's $\mu_T(E)$, paper.tex
  1108--1109, to distinguish it notationally from $\mathrm{massOfFunctional}(T)$ in the next step),
  this gives $\mathbb{M}_{\mathrm{AK}}(T) \le c$.

  \textbf{The min-vs-inf bridge, as an explicit step.} It remains to identify
  $\mathbb{M}_{\mathrm{AK}}(T)$ with $\mathrm{massOfFunctional}(T)$, since the two are defined
  differently: $\mathbb{M}_{\mathrm{AK}}(T) := \mu_T(E)$ for the specific \emph{minimum} admissible
  measure $\mu_T$ (paper.tex 1108, ``the minimum over all the finite positive Borel measure $\mu$
  satisfying'' the admissibility inequality \eqref{massmeasure}), while
  $\mathrm{massOfFunctional}(T)$ is by definition the \emph{infimum} of $\mu(E)$ over the class of
  admissible measures (\noderef{massOfFunctional}). We need only the free direction of this
  identification, and it costs nothing beyond existence-plus-admissibility of $\mu_T$: paper.tex
  1103--1108 asserts, for every $T \in \mathcal{M}_1(E)$ (in particular for our $T$, since $T :
  \mathrm{MetricCurrent1}\ E$), the existence of a finite positive Borel measure $\mu_T$ satisfying
  \eqref{massmeasure} -- i.e.\ $\mu_T$ is admissible for $T$ in the sense of
  \noderef{IsAdmissible} -- and paper.tex 1109 sets $\mathbb{M}_{\mathrm{AK}}(T) := \mu_T(E)$. Since
  $\mathrm{massOfFunctional}(T)$ is an infimum over the (nonempty, by the above) class of
  admissible measures, and $\mu_T(E)$ is the value of that same quantity $\mu(E)$ at one particular
  admissible measure $\mu = \mu_T$, the infimum is $\le$ that value:
  \[
    \mathrm{massOfFunctional}(T) \le \mu_T(E) = \mathbb{M}_{\mathrm{AK}}(T) \le c
    .
  \]
  (No attainment of the infimum and no minimality of $\mu_T$ among admissible measures is used
  here -- only that $\mu_T$ \emph{is} admissible, which is exactly what paper.tex 1103--1108
  quotes.) Converting the real inequality $\mathrm{massOfFunctional}(T) \le c$-bound into the
  $\overline{\mathbb{R}}_{\ge 0}$-valued conclusion via $c \ge 0$ gives $\mathrm{massOfFunctional}(T)
  \le \mathrm{ENNReal.ofReal}(c)$, as required.
\end{definition}
